% Generated from tex/chapters/ch04/08_構築技術.tex for the SWEBOK structure tree
\subsubsection{オブジェクト指向の実行時問題}

オブジェクト指向言語には、実行時に振る舞いを決める仕組みがある。代表例は多態性と反射である。

多態性とは、一般的な操作を、実際の具体的なオブジェクトの種類に応じて実行時に決められる性質である。たとえば、同じ「描画する」という操作でも、円、四角形、文字列で実行される処理が異なる。このように実行時に呼び出し先が決まることを動的束縛という。

反射とは、プログラムが実行時に自分自身の構造や振る舞いを調べたり、変更したりできる性質である。たとえば、クラス名やメソッド名を文字列として扱い、実行時にオブジェクトを生成したり、メソッドを呼び出したりできる。

これらは柔軟性を高めるが、理解やテストを難しくすることがある。利用する場合は、意図を明確にし、過度に動的な構造を避ける必要がある。
